\scanpage{386}
\mspage{110}
% archive sheet number: 385
et si de plus on suppose que l'on est dans l'un des deux cas suivants

a) $S$ est un ens. ordonné fini (ce qui nous amène à nous restreindre, à défaut de mieux, à $\mathrm{Diag}={}$\unsure{ss}-catégorie de $\mathrm{Cat}$ formée des ens. ord finis).

b) La factorisation stipulée dans l'axiome de factorisation $f=pi$ [avec $i\in\mathrm{Cof}$, $p\in W$ (cas des cofibrations) ou $i\in W$, $p\in\mathrm{Fib}$ (cas des fibrations) peut-être choisie fonctorielle en $f$.

Or, dans le cas de $\mathrm{Cat}$, dans les cas où on sait prouver que $W$ est associé à une structure à cofibrations (axiome $\underline{\mathrm{Loc}(6a)}$) ou à fibrations (axiome $\underline{\mathrm{Loc}(7\mathrm{D}_{W_{\unsure{\omega|w|\infty}}})}$), on trouve bel et bien une factorisation fonctorielle. Donc chacun de ces axiomes\unsure{,} \deleted{ainsi} implique $\mathrm{Der}\,2$\unsure{:|.}

